\paragraph{A geometric instance with every primitive arrow singular.}
The preceding distinction has an actual finite-field realization.
There are a finite extension $k/\mathbf F_{256}$, a smooth
projective curve $Y/k$, and a connected cyclic etale cover
$D'\to Y$ of degree thirteen such that every primitive
Frobenius character arrow is singular, but the actual primitive
factor has positive geometric two-rank. More precisely, one
can arrange
\[
 \begin{gathered}
 d=5\cdot3^a,\qquad 13\le a\le65,\qquad g(Y)=d+1,\\
 \dim P'_{13}=12d,\qquad 12\le f(P'_{13})\le12(d-1).
 \end{gathered}
 \tag{PIA.1}
\]
There is a geometric cyclic class $\mathcal M$ of exact order
thirteen on $Y^{(1)}$, defining the cover under relative
Frobenius, with
\[
 H^0\bigl(Y^{(1)},B_Y\otimes\mathcal M^i\bigr)\ne0
                           \qquad(1\le i\le12).
 \tag{PIA.2}
\]
Thus all primitive classes meet the actual locally exact
differential theta divisor, yet the primitive Frobenius operator
is not nilpotent. This is an existence corollary, not an
explicit equation or an exact computation of the new two-rank.

\begin{proof}
\emph{1. Verify the degree support in the cited construction.}
Work first over $\overline{\mathbf F}_2$ with the fixed genus-two
curve $C$ and cyclic thirteen-cover $D\to C$ of A32.
Its primitive factor $P_{13}$ is ordinary of dimension twelve,
so $f(P_{13})=12$.

For genus two and characteristic two, Raynaud's proof of
Theorem~2 makes the explicit choices $m=5$ and $n=3$
on printed p.~86. The first stage is a cyclic etale cover
$C'\to C$ of degree five; its genus is six.
Corollary~18 applies on $C'$ with the usual principal
polarization: $3\mid6$ and $6/3=2\le5=(6-1)/(2-1)$.
In the construction preceding Theorem~17, printed p.~84,
the resulting etale Galois cover $Z\to C'$ has group
\[
 1\longrightarrow\mu_3\longrightarrow H
            \longrightarrow J_{C'}[3]\longrightarrow1,
                         \qquad |H|=3^{13}.
 \tag{PIA.3}
\]
It supplies a representation $\rho'$ of $\pi_1(C')$ factoring
through $H$ with no theta divisor
\cite[pp.~84--86, Theorem~17, Corollary~18 and proof of Theorem~2]{RaynaudNonordinary2000}.

Take the Galois closure $Y\to C$ of $Z\to C$ geometrically.
Over $C'$ its function field is the compositum of the five
conjugates of the $H$-extension. Restriction embeds its Galois
group over $C'$ into $H^5$, and maps it onto $H$ for the
original conjugate. Hence that group has order $3^a$ with
$13\le a\le65$, and $[Y:C]=5\cdot3^a=d$.
In particular thirteen does not divide this degree.
This argument verifies the support after Galois closure;
it does not assume that the first Heisenberg cover was
already Galois over $C$.

\emph{2. Retain the theta-containment conclusion.}
The induced representation
$\rho=\operatorname{Ind}_{\pi_1(C')}^{\pi_1(C)}\rho'$
factors through the Galois group of $Y/C$: its kernel
contains the core of the subgroup defining $Z$.
By Raynaud's Proposition~1(4), a representation without a
theta divisor still has none after this induction.
Proposition~1(5) therefore gives, for $\alpha:Y\to C$,
\[
       \alpha^{(1)*}J_{C^{(1)}}
                         \subset\operatorname{Supp}(\Theta_Y)
 \tag{PIA.4}
\]
\cite[Proposition~1(4)--(5), p.~75]{RaynaudNonordinary2000}.

Choose the exact-order-thirteen line bundle $\mathcal L$
on $C^{(1)}$ defining the cyclic cover corresponding to
$D\to C$ under the relative-Frobenius equivalence of etale
covers, and put $\mathcal M=\alpha^{(1)*}\mathcal L$.
Norm composed with pullback is multiplication by $d$.
Since $13\nmid d$, $\mathcal M$ still has order thirteen.
Its relative-Frobenius pullback defines
$D'=D\times_CY\to Y$. Every $\mathcal M^i$, including all
twelve primitive powers, belongs to the image in (PIA.4).
The defining support property of the actual theta divisor
gives (PIA.2).

\emph{3. Positive stable rank survives.}
The covers $D/C$ and $Y/C$ have coprime geometric degrees
thirteen and $d$. They are linearly disjoint, so $D'$ is
geometrically connected and remains cyclic of degree thirteen
over $Y$. The actual-projector injection (PI.2) applies to
$D'\to D$, whose degree $d$ is odd. It gives
\[
                            f(P'_{13})\ge f(P_{13})=12.
 \tag{PIA.5}
\]
Riemann--Hurwitz gives $g(Y)=d+1$ and $g(D')=13d+1$;
since thirteen is prime, the primitive factor is the whole
Prym and has dimension $12d$.

For each $1\le i\le12$, the locally exact differential
sequence twisted by $\mathcal M^i$ identifies the kernel of
the corresponding relative-Frobenius map on coherent $H^1$
with $H^0(B_Y\otimes\mathcal M^i)$. The degree-zero line
bundles at either end have no global sections because they
are nontrivial of order thirteen. Each character space has
dimension $g(Y)-1=d$, and (PIA.2) makes every such map
singular. This is the character interpretation in
Pop--Sa\"idi, Section~1.3, printed p.~112 \cite{PopSaidi}.
The primitive coherent space is their direct sum, of dimension
$12d$; its Frobenius kernel has dimension at least twelve.
Its stable rank is consequently at most $12d-12$.
Together with (PIA.5) this proves (PIA.1), including the
failure of nilpotence despite every primitive arrow being
singular.

\emph{4. Descend to a finite field.}
All the curves, maps and the finitely many required deck
transformations over $\overline{\mathbf F}_2$ descend to
a common finite extension $k$ of $\mathbf F_{256}$.
Geometric connectedness, degrees, theta incidences and
geometric two-ranks are unchanged by this descent.
If a rational pointing is desired, a further finite extension
rationalizes a chosen point of $Y$ above A32's split point
and hence a point of the fibre product. No original-field
rational-unit assertion is made after such extensions.
\end{proof}

This corollary strengthens the linear-algebra illustration to
actual curves over a finite field. It does not supply a new
primitive-unit witness, and no main-gate advance is inferred
from it. The earlier A04 already refuted the weaker logical
inference; the present construction records a more specific
higher-genus failure with all primitive theta arrows present.
